\appendix
\setcounter{figure}{0} 
\renewcommand{\thesection}{Appendix~\Alph{chapter}.\arabic{section}}
\refstepcounter{chapter} 
\chapter*{Frequently Asked Questions} 
\markboth{Appendix~\thechapter\quad Frequently Asked Questions}%
         {Appendix~\thechapter\quad Frequently Asked Questions}
\addcontentsline{toc}{chapter} {\protect\numberline{Appendix}}
\addcontentsline{toc}{chapter} {\protect\numberline{\thechapter} Frequently Asked Questions}

\def\theHfigure{appendix.\arabic{figure}}

\noindent \textbf{Q1. When Someone You Love Passed Away, and You Cannot Find Your Way Through the Grief}

You must miss deeply the brother or sister with whom you shared so much of life's journey. And perhaps your longing for our heavenly home has grown stronger than ever. I believe the Lord knows the ache of your heart and understands how difficult this journey of separation can be. Yet the Lord has chosen to leave you here on earth for a time. You have not been left behind to bear your sorrow alone. You are caught up in grace, and one day will be brought into His glory.

One day, we will be with Christ, the Lord of heaven and earth. Our greatest hope is to dwell with Him forever and behold His glory face to face. 

Dear sister, as you walk through the days on earth, may you continue to fix your eyes on Christ. He is not only your refuge in this life but also your hope for eternity. May the Lord Himself sustain you, so that day by day you may experience His presence and grace.

The Lord has promised that He will wipe away our tears.
Ecclesiastes 3:1, 4 reminds us that there is a time for everything: a time to weep and a time to laugh, a time to mourn and a time to dance.

If you begin to experience joy again, it does not mean that you love your husband any less. If you travel or attend a gathering, it does not mean that you have forgotten him. You may continue to cherish the memories of the life you shared, even as you slowly learn to live with his absence. 

\vspace{0.3cm}

\noindent \textbf{Q2. At the Resurrection, People Will Neither Marry nor Be Given in Marriage; They Will Be Like the Angels in Heaven (Matthew 22:30)}

I can only imagine the scene in Revelation 7:9--10:

\begin{quote}
``After this I looked, and there before me was a great multitude that no one could count, from every nation, tribe, people and language, standing before the throne and before the Lamb. They were wearing white robes and were holding palm branches in their hands. And they cried out in a loud voice:

`Salvation belongs to our God, who sits on the throne, and to the Lamb.' ''
\end{quote}

Perhaps, amid that vast multitude, you will catch sight of your husband.

Yet the one standing before you will no longer appear just as you knew him on earth. Clothed in white, his eyes shining with joy, he will be praising the greatness of salvation. With his whole being, he will worship the Lord who redeemed him. His life will be centered entirely on Christ. His deepest desire will be to draw near to the Lord, to know Him more fully, to praise Him, and to walk with Him throughout eternity.

You may discover that the man you meet there is far more glorious than the man you knew in this life.

You knew him well. You knew how he devoted himself in his work. You knew that he used to drink a cup of coffee before leaving for work each morning. You were familiar with his joys and sorrows, his laughter and frustration. You walked beside him through life's rises and falls. Yet on that day, he will no longer be constrained by the weakness, toil, and sorrow of this present life. He will have received an imperishable, glorified body (1 Corinthians 15:42--44).

But you, too, will have changed.

You will be standing among that same multitude from every nation, tribe, people, and language, clothed in white and holding a palm branch in your hand. Together, you will offer praise to God, who sits on the throne, and to the Lamb. Before you will be the Savior who knows you more completely than anyone else. You will love and worship Him and walk with Him forever.

As 1 Corinthians 13:12 says, ``For now we see only a reflection as in a mirror; then we shall see face to face.'' Our knowledge of the Lord is limited in this life, but one day we will know Him more fully. 

\vspace{0.3cm}


\noindent \textbf{Q3. If I Am in Heaven One Day, but the Family Members I Love Are Not There, Will I Still Be Able to Rejoice?}

When we think of our parents, spouse, children, or siblings, it is difficult to imagine how we could rejoice in heaven if those we love are not there. Will we remember them? If we do, how can we be free from sorrow?

Scripture does not fully explain how the redeemed will remember or understand the final destiny of loved ones who are not saved. We must therefore approach this question with humility, trusting what God has revealed without speculating beyond His Word.

\textbf{First}, will we remember those we love?

We know that after His resurrection, Jesus remained the same Lord who had been crucified and then raised from the dead. In a similar way, God will transform our lowly bodies so that they will be like His glorious body (Philippians 3:21). Yet Scripture does not explicitly tell us how our memories of those who were not saved will remain, or how we will understand those memories in eternity. We should be careful not to claim more than God has revealed.

\textbf{Second}, how can we rejoice if someone we love is not in heaven?
Revelation 21:4 promises:

\begin{quote}
``He will wipe every tear from their eyes. There will be no more death or mourning or crying or pain, for the old order of things has passed away.''
\end{quote}

God promises that sorrow and pain will be no more. Although Scripture does not explain how this promise relates to our memories of those who are lost, we can trust that God's promise will never fail. Revelation 16:7 declares, ``Yes, Lord God Almighty, true and just are your judgments.'' We can trust both His perfect justice and His unfailing goodness, even when we do not understand how He will resolve this mystery.

\textbf{Third}, God's holiness and compassion do not contradict one another.
Jesus knew the reality of sin and judgment, yet He grieved over Jerusalem (Matthew 23:37). His example reminds us that compassion for sinners is not incompatible with God's holiness and justice.

How, then, should we respond today? Paul expressed ``great sorrow and unceasing anguish'' for his fellow Israelites (Romans 9:2–3). His example reminds us that grieving over the spiritual condition of those we love is not evidence of weak faith.

We should continue to pray for our loved ones, share the gospel, and trust them to God. We do not know every person's final destiny, nor can we fully understand the mysteries of eternity. But we know the character of the God who holds eternity in His hands. He is perfectly just, unfailingly good, and faithful to every promise He has made.
